\documentclass[UTF8,a4paper,11pt]{ctexart}
\usepackage[margin=2.5cm]{geometry}
\usepackage{booktabs,tabularx,url}
\usepackage[hidelinks]{hyperref}
\setlength{\parindent}{2em}
\title{AI工具使用详情}
\author{}
\date{2026年9月10日}
\begin{document}
\maketitle

\section{工具与使用范围}
本任务使用OpenAI Codex桌面应用及其辅助代理、代码执行、文件处理和文献检索功能。模型名称以实际会话或客户端显示为准，本文档不编造未核验的内部子版本。先前ChatGPT对话“母模型建模思路”和“论文收集与匹配”作为建模路线及检索线索使用；其中数学表达和书目内容并未直接视为事实，而是与赛题、原始附件及可获得的一手文献信息核对。

AI参与了统一建模方案整理、预测与优化程序编写、数据计算、场景与消融实验设计、图表生成、书目信息核验、论文初稿撰写和附件5数据填报。其参与范围包含实质性的模型及代码内容，并不限于语言润色。

\section{主要提示与过程}
用户要求结合上述两段对话，遵循C题，在指定LaTeX模板目录完成论文并填写附件5，保留原始代码，同时明确由用户后续编译与审查。过程中用户进一步要求保留附件5原始时间标签，从0:10--0:20开始，以次日0:00--0:10结束。上述约束被纳入数据导出及论文的时间映射说明。

工作过程依次包括：读取赛题与模板，提取附件1至4；按可获得信息选择历史预测，保留日内相关残差；建立线性规划与当前时段反馈执行器；从1月初态连续推进全年；生成指定日期、全年策略和对照实验；依据真实计算结果撰写正文，输出原模板工作簿与源代码。原始附件及模板代码备份另存，未用伪造实测结果填表。

\section{已实施的计算核对}
程序核对了数据维度、缺失与非负性、部分预测的未来信息扰动不变性、逐时供需平衡、储电量边界、功率上限、充放电互斥、跨日状态连续、费用分解，以及结果工作簿与求解数组的一致性。文献书目使用DOI注册信息及出版社入口核验，并明确了摘要或可用全文不能支持的细节。以上核对由计算程序及AI协助完成，不等于参赛人员已经逐项人工确认。

\section{采纳、人工修改与核验状态}
本次交付内容为AI协助形成的论文源码、计算源程序与结果工作簿。按用户要求，本次没有安装LaTeX、没有编译论文，也没有生成论文审稿报告。用户尚未完成或未在本会话中记录的人工修改和人工核验，不在此声称已完成。

用户在最终采用前，应根据实际完成情况补充：参与人工核验人员及日期、模型和结算假设是否采纳、时间模板映射是否确认、主要代码与数值的人工复核记录，以及编译排版后的修改。补充内容应反映真实操作，不能将计划进行的核验写成已经完成。

\section{对应材料}
论文入口为\texttt{main.tex}；全部计算源文件位于\texttt{code/}；原模板位于\texttt{originals/}；原始数据摘要、实验设置与结果位于\texttt{data/}；五个结果工作簿保存于题目指定的附件5目录。本文件由用户后续单独编译为AI工具使用详情PDF。
\end{document}
